% GENERATED FILE. Edit canonical lessons, then run npm run generate:books.
% canonical-source-hash: fnv1a64:340dec2b9c5735b1
% canonical-lessons: TE-C78-isit, TE-C78-doyouknow, TE-C78-dontknow, TE-C78-didntunderstand, TE-C78-isthere

\chapter{Yes or No, and I Do Not Know}
\label{ch:te-yes-or-no}
\hlchaptermodality{78}

\begin{chapteropening}
\textbf{By the end of this chapter:} \emph{I can ask a Telugu yes-or-no question, say I do not know, say I did not understand, and ask whether somebody is there.}

Complete TE-C78-isthere and demonstrate the chapter promise.
\end{chapteropening}

\section[avunā]{\te{అవునా} (avunā) — is it so?}

\label{lesson:TE-C78-isit}



\begin{quote}
Before the new one: what did \emph{muppai} mean?
\end{quote}

\subsection*{You'll want to know: \te{అవునా}}
\textbf{\te{అవునా}} (\emph{avunā}) — \textquotedblleft{}is it so?\textquotedblright{}.

Every question you have asked so far needed a question \emph{word}. This one needs none: take \te{అవును}, \textquotedblleft{}yes\textquotedblright{}, and add \textbf{-\te{ఆ}}.

That \textbf{-\te{ఆ}} is the whole yes-or-no machine. Hang it on the end of any statement and it becomes a question.

\begin{grammarlens}[title={Grammar Lens: -\te{ఆ} makes any statement a question}]
\textbf{\te{మీరు} \te{వెళ్తారు}} \textquotedblleft{}you go\textquotedblright{} becomes \textbf{\te{మీరు} \te{వెళ్తారా}?} \textquotedblleft{}do you go?\textquotedblright{}. \textbf{\te{తెలుగు} \te{కష్టం}} becomes \textbf{\te{తెలుగు} \te{కష్టమా}?}. Nothing moves and no word is added: the last vowel lengthens into \te{ఆ} and the sentence is a question. This is the commonest question type in any listening paper, and it is one ending.
\end{grammarlens}

\begin{grammarlens}[title={What you've built}]
A question with no question word in it.
\end{grammarlens}

\subsection*{Your turn}
\begin{itemize}
\raggedright
  \item \emph{Say it:} \emph{avunā}
  \item \emph{Say it:} \emph{avunā}, once more
  \item \emph{Say it:} say \emph{avunu}, then \emph{avunā}, and hear the difference
  \item \emph{Recall:} say \emph{mancidi}, then \emph{mukku}, then say \emph{avunā}
\end{itemize}

\begin{tcolorbox}[breakable,colback=teal!4,colframe=teal!35!black,title={Before you move on}]
What does \te{అవునా} mean? (\textquotedblleft{}is it so?\textquotedblright{}.) Say it once more.
\end{tcolorbox}
\section[telusā]{\te{తెలుసా} (telusā) — do you know?}

\label{lesson:TE-C78-doyouknow}



\begin{quote}
Before the new one: what did \emph{avunā} mean?
\end{quote}

\subsection*{You'll want to know: \te{తెలుసా}}
\textbf{\te{తెలుసా}} (\emph{telusā}) — \textquotedblleft{}do you know?\textquotedblright{}.

The first thing you will do with the new ending. You already have \textbf{\te{తెలుసు}}, \textquotedblleft{}it is known\textquotedblright{}; add \textbf{-\te{ఆ}} and you are asking.

\textbf{\te{మీకు} \te{తెలుగు} \te{తెలుసా}?} — \textquotedblleft{}do you know Telugu?\textquotedblright{} The knower sits in the dative \te{మీకు}, the way it did when you first met \te{తెలుసు}.

\begin{grammarlens}[title={What you've built}]
A real question, built from a word you already had.
\end{grammarlens}

\subsection*{Your turn}
\begin{itemize}
\raggedright
  \item \emph{Say it:} \emph{telusā}
  \item \emph{Say it:} \emph{telusā}, once more
  \item \emph{Say it:} ask \emph{mīku telugu telusā?}
  \item \emph{Recall:} say \emph{sulabhaṁ}, then \emph{nōru}, then say \emph{telusā}
\end{itemize}

\begin{tcolorbox}[breakable,colback=teal!4,colframe=teal!35!black,title={Before you move on}]
What does \te{తెలుసా} mean? (\textquotedblleft{}do you know?\textquotedblright{}.) Say it once more.
\end{tcolorbox}
\section[nāku teliyadu]{\te{నాకు} \te{తెలియదు} (nāku teliyadu) — I do not know}

\label{lesson:TE-C78-dontknow}



\begin{quote}
Before the new one: what did \emph{telusā} mean?
\end{quote}

\subsection*{You'll want to know: \te{నాకు} \te{తెలియదు}}
\textbf{\te{తెలియదు}} (\emph{teliyadu}) — \textquotedblleft{}it is not known\textquotedblright{}, which is Telugu for \textquotedblleft{}I don't know\textquotedblright{}.

This is the answer to the last word, and it is probably the sentence you will need most often of anything in this book.

It is not \te{తెలుసు} with \te{లేదు} after it. Telugu negates an ordinary verb from the inside, with \textbf{-\te{అదు}} on the stem: \te{తెలియ}- + -\te{దు}.

\begin{grammarlens}[title={Grammar Lens: how Telugu says a verb did not happen}]
The two negatives of your first lesson, \te{లేదు} and \te{కాదు}, deny \textbf{existence} and \textbf{identity}. An ordinary action is denied differently — with an ending on the verb itself: \te{వెళ్తాను} \textquotedblleft{}I go\textquotedblright{} becomes \textbf{\te{వెళ్ళను}} \textquotedblleft{}I don't go\textquotedblright{}; \te{వస్తాను} becomes \textbf{\te{రాను}}. The ending carries the negative and no separate word is used. \te{తెలియదు} is that pattern on \te{తెలియ}-.
\end{grammarlens}

\begin{grammarlens}[title={What you've built}]
The most useful sentence in the chapter.
\end{grammarlens}

\subsection*{Your turn}
\begin{itemize}
\raggedright
  \item \emph{Say it:} \emph{teliyadu}
  \item \emph{Say it:} \emph{teliyadu}, once more
  \item \emph{Say it:} answer \emph{mīku telugu telusā?} with \emph{teliyadu}
  \item \emph{Recall:} say \emph{kaṣṭaṁ}, then \emph{guṇḍe}, then say \emph{nāku teliyadu}
\end{itemize}

\begin{tcolorbox}[breakable,colback=teal!4,colframe=teal!35!black,title={Before you move on}]
What does \te{నాకు} \te{తెలియదు} mean? (\textquotedblleft{}I do not know\textquotedblright{}.) Say it once more.
\end{tcolorbox}
\section[arthaṁ kālēdu]{\te{అర్థం} \te{కాలేదు} (arthaṁ kālēdu) — I did not understand}

\label{lesson:TE-C78-didntunderstand}



\begin{quote}
Before the new one: what did \emph{teliyadu} mean?
\end{quote}

\subsection*{You'll want to know: \te{అర్థం} \te{కాలేదు}}
\textbf{\te{అర్థం} \te{కాలేదు}} (\emph{arthaṁ kālēdu}) — \textquotedblleft{}I did not understand\textquotedblright{}.

You already have \textbf{\te{నాకు} \te{అర్థమైంది}}, \textquotedblleft{}I understood\textquotedblright{} — literally \textquotedblleft{}to me it became meaning\textquotedblright{}. This is the same sentence with the becoming denied: \te{కా}- \textquotedblleft{}become\textquotedblright{} plus \te{లేదు}.

Say it, and the person you are talking to will slow down. That is the whole point of a repair phrase.

\begin{grammarlens}[title={What you've built}]
A way out of the moment you lose the thread.
\end{grammarlens}

\subsection*{Your turn}
\begin{itemize}
\raggedright
  \item \emph{Say it:} \emph{arthaṁ kālēdu}
  \item \emph{Say it:} \emph{arthaṁ kālēdu}, once more
  \item \emph{Say it:} say \emph{nāku arthamaindi}, then \emph{arthaṁ kālēdu}
  \item \emph{Recall:} say \emph{nunci}, then \emph{hṛdayaṁ}, then say \emph{arthaṁ kālēdu}
\end{itemize}

\begin{tcolorbox}[breakable,colback=teal!4,colframe=teal!35!black,title={Before you move on}]
What does \te{అర్థం} \te{కాలేదు} mean? (\textquotedblleft{}I did not understand\textquotedblright{}.) Say it once more.
\end{tcolorbox}
\section[unnārā]{\te{ఉన్నారా} (unnārā) — is he or she there?}

\label{lesson:TE-C78-isthere}



\begin{quote}
Before the new one: what did \emph{arthaṁ kālēdu} mean?
\end{quote}

\subsection*{You'll want to know: \te{ఉన్నారా}}
\textbf{\te{ఉన్నారా}} (\emph{unnārā}) — \textquotedblleft{}is he or she there?\textquotedblright{}.

\te{ఉండు} in its respectful \textbf{-\te{ఆరు}} form, with the new \textbf{-\te{ఆ}} on the end. Three pieces you already had.

It is what you say at a door or into a telephone: \textbf{\te{గురువు} \te{ఉన్నారా}?} — \textquotedblleft{}is the teacher there?\textquotedblright{}

\begin{grammarlens}[title={What you've built}]
A question you can ask at somebody's door.
\end{grammarlens}

\subsection*{Your turn}
\begin{itemize}
\raggedright
  \item \emph{Say it:} \emph{unnārā}
  \item \emph{Say it:} \emph{unnārā}, once more
  \item \emph{Say it:} ask whether the teacher is there
  \item \emph{Recall:} say \emph{paṭṭaṇaṁ}, then \emph{pālu}, then say \emph{unnārā}
\end{itemize}

\begin{tcolorbox}[breakable,colback=teal!4,colframe=teal!35!black,title={Before you move on}]
What does \te{ఉన్నారా} mean? (\textquotedblleft{}is he or she there?\textquotedblright{}.) Say it once more.
\end{tcolorbox}
